\documentclass[pdflatex,sn-basic,Numbered]{sn-jnl}% numbered references, SN Computer Science

%% Packages (sections are generated from ../paper by make_sections.py)
\usepackage{graphicx}
\usepackage{multirow}
\usepackage{amsmath,amssymb,amsfonts}
\usepackage{amsthm}
\usepackage{mathrsfs}
\usepackage{booktabs}
\usepackage{algorithm}
\usepackage{algorithmic}
\usepackage{bm}
\usepackage{enumitem}
\usepackage{placeins}
\usepackage{xcolor}
\usepackage{textcomp}
\usepackage{manyfoot}
\usepackage{varwidth}
\def\UrlBreaks{\do\/\do\-\do\.\do\_\do\?\do\=\do\&}
\Urlmuskip=0mu plus 1mu

\theoremstyle{thmstyleone}
\newtheorem{theorem}{Theorem}
\newtheorem{proposition}{Proposition}
\newtheorem{corollary}{Corollary}
\newtheorem{lemma}{Lemma}
\theoremstyle{thmstyletwo}
\newtheorem{example}{Example}
\newtheorem{remark}{Remark}
\theoremstyle{thmstylethree}
\newtheorem{definition}{Definition}

\raggedbottom

\begin{document}

\title[Evasion of Tree Ensembles in Network Intrusion Detection]{Steering Through the Trees: Physical Lie Vector Field Evasion of Tree Ensembles in Network Intrusion Detection}

\author*[1]{\fnm{Quang-Vinh} \sur{Dang}}\email{vinh.dq4@buv.edu.vn}

\affil*[1]{\orgdiv{School of Innovative and Computing Sciences}, \orgname{British University Vietnam}, \orgaddress{\city{Hung Yen}, \country{Vietnam}}}

\abstract{\textbf{Purpose:} Gradient-boosted decision trees dominate flow-based network intrusion detection, yet their worst-case robustness to schema-feasible evasion is rarely measured. We ask how far an adversary with score access only can evade them, and whether a defense can bound what remains.
\textbf{Methods:} Flow-Field Evasion (FFE) models packet-level actions as commuting generators on the flow manifold. After diagnosing FFE-v1, we introduce FFE-v2 (re-segmentation, multi-scale probing, data-driven windows), which uses scores only. We compare it with FFE-v1, constrained adaptive projected gradient descent (CAPGD) and our own greedy white-box control on three full, deduplicated corpora and three tree libraries, and add a defense that OR-fuses a monotone gate with a free classifier.
\textbf{Results:} FFE-v2 matches or exceeds the control and FFE-v1 in all nine cells (significantly above the control in four). On XGBoost its real attack success rate (ASR) is 99.5\%, 100.0\% and 98.3\%, against 57.6\% for FFE-v1 and 21.1\% for CAPGD on UNSW-NB15, with 28--38 queries. CatBoost on CIC-IDS2017 (36.4\%) is a ceiling set by immutable features. Against FFE-v1, CAPGD and a greedy white-box attacker, the defense leaves a real ASR of at most 8.8\% on CSE-CIC-IDS2018 and CIC-IDS2017 and 1.0--32\% on UNSW-NB15.
\textbf{Conclusion:} Feasibility is relative to a schema verifier, and chaff accounting matters on UNSW-NB15. Flagged flows stay flagged only if backward traffic is unchanged; hiding those features costs accuracy and weakens the defense. FFE-v2 was run against it only on CIC-IDS2017 XGBoost. The single test evaluation is not a clean hold-out; five-seed runs are cleaner.}

\keywords{Adversarial machine learning, Network intrusion detection, Gradient boosted decision trees, Evasion attacks, Problem-space constraints, Lie algebra}

\maketitle
\sloppy

\section{Introduction}
\label{sec:intro}

Modern network intrusion detection systems (NIDS) increasingly rely on machine learning over flow-based telemetry. At multi-gigabit rates, inspecting every packet payload is prohibitive, so sensors such as CICFlowMeter \cite{sharafaldin2018toward} aggregate bidirectional packet exchanges into tabular flow records of timing, length, header and flag statistics. Gradient-boosted decision trees (GBDTs), notably XGBoost \cite{chen2016xgboost}, LightGBM \cite{ke2017lightgbm} and CatBoost \cite{prokhorenkova2018catboost}, are the dominant classifiers for these records \cite{dang2022using, dang2022studying}.

Near-perfect benchmark accuracy does not guarantee robustness: attackers control the malleable parts of their traffic, such as timing, padding, segmentation and window sizes, so a detector that relies on these features can be evaded without losing the exploit \cite{apruzzese2022modeling, sommer2010outside, geirhos2020shortcut}.

Evading a tree ensemble is hard for two reasons. First, the ensemble is piecewise constant, so $\nabla_x F(x)=0$ for almost all $x$ and gradient-based attacks receive no signal. Second, flow features obey conservation laws. Flow Bytes/s satisfies $x_8 = 10^6(x_4+x_5)/x_6$, where $x_4$ and $x_5$ are the forward and backward byte totals and $x_6$ is the duration in microseconds. A perturbation that changes the byte totals without updating the rate produces a phantom flow that protocol sanity checks reject \cite{pierazzi2020intrusion, apruzzese2022modeling}. Existing attacks either ignore this coupling, as ZOO \cite{chen2017zoo} and HopSkipJump \cite{chen2020hopskipjump} do (HopSkipJump yields no feasible evasive flow on the corpora of this study; ZOO keeps at most 25.2\% Real ASR, the share of attacks that evade and satisfy the schema invariants), or need hundreds of black-box queries per flow \cite{simonetto2024constrained, ghamizi2020moeva, han2021traffic}, or depend on a surrogate \cite{simonetto2024constrained, zolbayar2022nidsgan}.

Our earlier attack, Flow-Field Evasion (FFE-v1), derived four packet-level actions as closed-form generators of a Lie cone on the flow manifold, so that every step preserves the invariants, and steered with score probes. It evaded XGBoost in 99.5\%, 99.2\% and 57.6\% of the detected attacks of three full corpora, but a rigorous control exposed its limits: a greedy white-box attack in the style of leaf-tuple search \cite{zhang2020efficient}, which reads the trees but projects every candidate onto the manifold, reached 90.2\% on UNSW-NB15, and FFE-v1 evaded LightGBM on CIC-IDS2017 in only 35.1\% and CatBoost on CSE-CIC-IDS2018 in only 59.3\% of the flows. We diagnosed these gaps on a validation split, not on the test cohorts: the attack lacked a generator that raises the packet count without adding payload, used one fixed magnitude per generator and one hard-coded window target (64,240~B) that lies outside the support of the CIC-IDS2017 data, and measured progress on a saturating probability.

This article reports the repaired attack, FFE-v2, and a defense built on what the attack revealed. The main contributions are as follows.
\begin{itemize}
    \item \textbf{Generators.} We extend the Lie-cone model by a re-segmentation generator (the same payload carried by more, smaller packets). The five fields still commute and preserve the rate, mean-length, header and MTU invariants of the schema, although they are not linearly independent (re-segmentation lies in the span of chaff and padding, outside their nonnegative cone; Section~\ref{sec:framework}, Online Resource~1, Section~S1).
    \item \textbf{Attack.} FFE-v2 probes every generator on a geometric ladder of magnitudes, draws window targets from the attacker's benign sample, measures progress in log-odds and refines the cheapest crossing; it needs only scores (Section~\ref{sec:algorithm}). In a single evaluation on frozen settings it matches or exceeds our greedy threshold-crossing white-box control in all nine corpus-library cells (significantly above it in four) and is significantly above FFE-v1 in seven (Real ASR 99.5\%, 100.0\% and 98.3\% on XGBoost against 99.5\%, 99.2\% and 57.6\% for FFE-v1 and 88.5\%, 89.5\% and 21.1\% for CAPGD); five-seed means agree (Section~\ref{sec:results}). The control is our own implementation (erratic across seeds on CIC-IDS2017 LightGBM); LT-Attack and Cascioli et al.\ were not run.
    \item \textbf{Honest accounting.} FFE-v2 uses 28--38 score queries (109 on CIC-IDS2017 CatBoost), against 2 for the control and about 3 for CAPGD with an uncounted surrogate, and more than FFE-v1 on the easy corpora; its operational cost on UNSW-NB15 exceeds that of FFE-v1; one cell (CatBoost on CIC-IDS2017, 36.4\%) is a ceiling set by attacker-immutable features rather than an attack success. ``Feasible'' is relative to a schema verifier that does not check the header-plus-payload bound, and chaff probes are booked as $L_0$ bytes of the length column; a synthetic replay shows that UNSW-NB15 evasions depend on chaff accounting (Section~\ref{sec:ablation}).
    \item \textbf{Defense.} We OR-fuse a monotone gate, a tree ensemble that is non-decreasing in forward packets, forward bytes, duration and header bytes and does not read the window, rates or mean length, with a free classifier (augmentation-hardened on five of nine cells; the undefended model on the others). If the attacker leaves the backward features and the port and protocol unchanged (an assumption on which the guarantee rests), a flow the gate flags stays flagged (Section~\ref{sec:defense}, a certificate in the spirit of monotone classification \cite{incer2018adversarially}). Against FFE-v1 ($M=5,20$), CAPGD and a greedy white-box threshold-crossing attacker, Real ASR on the flows still detected falls to at most 8.8\% on CSE-CIC-IDS2018 and CIC-IDS2017 and to 1.0--32\% on UNSW-NB15, with validation accuracy, F1 and recall not falling by more than 0.5 point. We state where this fails: about 8\% of CSE-CIC-IDS2018 flows are not flagged by the gate and remain evadable, the stand-alone gate costs accuracy, hiding the backward features from the gate to remove the fixed-backward assumption costs much accuracy and weakens V2, and FFE-v2 was run against the hardened models on one cell only.
\end{itemize}
Caveat: FFE-v2 was tuned on validation cohorts and evaluated once on the test cohorts, but the FFE-v1 failures that motivated it were first seen on those test cohorts, re-segmentation was motivated by the control's moves on the same corpora, and the window pool uses the attacker's benign training sample (tuning queries are not counted); the five-seed study (retrained models, fresh cohorts, no tuning) is the cleaner confirmation. The attacks are red teaming that estimates what a schema-constrained adversary with score access can reach; no real network was attacked.

\section{Related Work and Problem Formulation}
\label{sec:sota}

\subsection{Related Work and State-of-the-Art Taxonomy}
\label{sec:related_work}

Early adversarial attacks on NIDS were transplanted from image classification \cite{szegedy2014intriguing, goodfellow2015explaining, papernot2016limitations, carlini2017towards}. Applied to flow features they reported high success rates, but the perturbed vectors were not network flows: they assign fractional values to packet counts, modify immutable fields such as destination port, and break the conservation between byte totals and packet counts. The problem-space view corrects this. Pierazzi et al.~\cite{pierazzi2020intrusion} showed that an adversarial perturbation is meaningful only when it is realized as a transformation of the physical object, such as a binary or a packet sequence, and that the transformation must preserve functionality and survive the feature extractor. Surveys and audits~\cite{apruzzese2023real, he2023adversarial, apruzzese2022modeling} argue that success rates overstate the risk once realistic constraints are enforced. Constrained tabular attacks address feasibility only in part: CAA \cite{simonetto2024constrained} combines a gradient-based attack (CAPGD) with the MOEVA genetic search \cite{ghamizi2020moeva}, which handles constraints by penalties, and Traffic Manipulator \cite{han2021traffic} mutates PCAP files with particle swarm search over intervals, padding and chaff.

Recent work makes the realism requirement concrete for NIDS. Merzouk et al.~\cite{merzouk2022practicality} show that many published evasion attacks cannot be realized on live traffic, Sheatsley et al.~\cite{sheatsley2022adversarial}, Debicha et al.~\cite{debicha2023advbot} and Zhang et al.~\cite{zhang2025explainable} craft domain-constrained, botnet or transferable flows, and El Shehaby et al.~\cite{elshehaby2026perturbability} score which flow features an attacker can perturb, the mutable/immutable split that FFE builds into its schema. Learned generators~\cite{lin2022idsgan, zolbayar2022nidsgan} trade query cost for a training phase, and tabular benchmarks~\cite{kireev2023cost, simonetto2024tabularbench, he2026tabattackbench} stress the role of constraints in reported success rates. The tree-ensemble literature moved toward faster and deployable attacks~\cite{cascioli2024faster}, evasion detection~\cite{devos2023detecting} and robust training~\cite{vos2022robust, lindqvist2025adversarial}. None combines closed-form coupled-invariant generators with score-only steering and a certified defense for the same threat model.

Exact methods exploit the piecewise-constant structure of trees. Kantchelian et al.~\cite{kantchelian2016evasion} cast minimal evasion against tree ensembles as a mixed-integer linear program. The formulation is exact but NP-hard and accepts only linear constraints, so it cannot express the nonlinear rate relations of flow records. MaxSAT encodings of tree ensembles \cite{ignatiev2022formal} and SAT-based explanations of decision trees \cite{izza2022tackling} provide formal explanations and robustness checks, and LT-Attack \cite{zhang2020efficient} attacks tree ensembles through a discrete search over leaf tuples. These methods operate in unconstrained feature space and do not enforce coupled network invariants unless a projection is added (Online Resource~1, Section~S7).

On the defensive side, robust training of tree ensembles \cite{chen2019robust, andriushchenko2019provably, calzavara2020treant, vos2021efficient} and adversarial retraining of neural NIDS \cite{hashemi2020enhancing} improve robustness within a feature-space threat model, and verification of tree ensembles can certify robustness for restricted structures \cite{calzavara2023verifiable, calzavara2025verifiableboosted}. The closest idea to our certificate is monotonic classification for malware detection \cite{incer2018adversarially}: if the attacker can only add content, a score that is monotone in the added features cannot be lowered. We transfer this argument to flow records, where the difficulty is the coupling between features, which we handle by hiding the derived columns from the certified component (Section~\ref{sec:defense}).

Online Resource~1, Table~S21 summarizes these paradigms. Methods that preserve physical feasibility either need hundreds of black-box queries per flow or depend on a surrogate; FFE steers along packet-level generators and chooses the direction with score probes, and the white-box control of Online Resource~1, Section~S7 shows the same moves are reachable with tree access.

\subsection{Valid Packet Streams, Flow Telemetry and the Twelve-Feature Schema}
\label{sec:formulation}
A packet stream $P=((p_1,t_1),\dots,(p_m,t_m))$ is an ordered sequence of packets with $0\le t_1\le\dots\le t_m$. Only streams in the admissible set $\mathcal{P}_{\text{valid}}$ are considered: every packet carries valid IPv4, TCP and UDP checksums \cite{rfc791,rfc9293,rfc768} and obeys the interface MTU, $\text{Len}_{\text{IP}}(p_i)\le 1500$ bytes. A flow is the bidirectional set of packets sharing a 5-tuple, split into forward (client-to-server) and backward packets, and a flow exporter such as CICFlowMeter \cite{sharafaldin2018toward} computes a deterministic, many-to-one mapping $\Phi:\mathcal{P}_{\text{valid}}\to\mathbb{R}^d$, $x=\Phi(P)$.

We use a standardized 12-dimensional schema of the core metrics that state-of-the-art models analyze on CSE-CIC-IDS2018 \cite{sharafaldin2018toward}: $x_0$, destination port; $x_1$, protocol (6 for TCP, 17 for UDP); $x_2$, forward packets; $x_3$, backward packets; $x_4$, forward length total (the exporter's Total Length of Forward Packets); $x_5$, backward length total; $x_6=(t_m-t_1)\times 10^6$, duration in microseconds; $x_7$, initial forward TCP window ($0\le x_7\le 65535$, $x_7\equiv 0$ for UDP); $x_8$, flow bytes per second; $x_9$, flow packets per second; $x_{10}$, mean forward packet length; and $x_{11}$, forward transport header bytes, at least the per-packet minimum header size times $x_2$. The coupled columns are fixed by
\begin{equation}
    x_8=\frac{10^6\,(x_4+x_5)}{x_6},\qquad x_9=\frac{10^6\,(x_2+x_3)}{x_6},\qquad x_{10}=\frac{x_4}{\max(1,x_2)}.
\end{equation}
The generators of Section~\ref{sec:framework} book a chaff probe as $L_0$ bytes of $x_4$ (the wire size of an empty segment), whereas a re-segmented packet moves existing bytes and leaves $x_4$ unchanged. This accounting is a modelling choice that matters empirically (Section~\ref{sec:ablation}).

\subsection{Feature Invariants and the Flow Manifold}
The indices $\mathcal{D}=\{0,\dots,11\}$ fall into constraint classes derived from network semantics. $\mathcal{I}_{\text{imm}}=\{0,1,3,5\}$ (port, protocol, backward packets and bytes) are fixed by the service and the server. $\mathcal{I}_{\text{mono}}^+=\{2,4,6,11\}$ can only grow: an attacker may add packets, pad payloads or delay transmissions, but cannot remove what the exploit needs. The window obeys $0\le x_7\le 65535\cdot\mathbb{I}_{x_1=6}$, and $\mathcal{I}_{\text{dep}}=\{8,9,10\}$ are the deterministic functions above. In addition, $x_{11}\ge x_2(20\,\mathbb{I}_{x_1=6}+8\,\mathbb{I}_{x_1=17})$ and $0\le x_4\le 1500x_2$.

The physical flow space is the discrete image of valid packet streams, $\mathcal{M}_{\text{flow}}=\{x\mid\exists P\in\mathcal{P}_{\text{valid}},\ \Phi(P)=x\}$, with integer packet counts, byte totals, header totals, window and duration (in microseconds, at most $1.1\times10^8\,\mu$s, below the 120-second exporter timeout). For continuous differential geometry we relax counts and duration to reals and obtain the \emph{ambient flow manifold}
\begin{gather}
    \widetilde{\mathcal{M}}_{\text{flow}}=\Big\{x\in\mathbb{R}^{12}\;\Big|\;x_0\in[1,65535],\ x_1\in\{6,17\},\ x_2\ge1,\ x_3,x_5\ge0,\notag\\
    0<x_6\le1.1\times10^8,\ H_{\min}x_2\le x_{11}\le H_{\max}x_2,\ 0\le x_4\le1500x_2,\notag\\
    0\le x_7\le65535\,\mathbb{I}_{x_1=6},\ x_{10}=\tfrac{x_4}{x_2},\ x_8=\tfrac{10^6(x_4+x_5)}{x_6},\ x_9=\tfrac{10^6(x_2+x_3)}{x_6}\Big\}
\end{gather}
with $H_{\min}=20\,\mathbb{I}_{x_1=6}+8\,\mathbb{I}_{x_1=17}$ and $H_{\max}=60\,\mathbb{I}_{x_1=6}+8\,\mathbb{I}_{x_1=17}$. Continuous flows are formulated on $\widetilde{\mathcal{M}}_{\text{flow}}$, and discrete execution maps trajectories back to $\mathcal{M}_{\text{flow}}$ through the closed-form projection $\mathcal{P}_{\mathcal{M}}$ (Section~\ref{sec:algorithm}).

\subsection{Threat Model}
The adversary starts from a malicious flow $x_{\text{orig}}\in\mathcal{M}_{\text{flow}}$ with $F(x_{\text{orig}})\ge0$ (probability of the malicious class at least 0.5) and seeks $x^*\in\mathcal{M}_{\text{flow}}$ with $F(x^*)<0$. Two access levels are used. The \emph{score-only} attacker queries the malicious probability of the target and nothing else; this is the headline setting of FFE-v2. The \emph{threshold} attacker additionally reads the split thresholds of the trees, which enables the optional snapping step (Section~\ref{sec:algorithm}) and is also the setting of the projected white-box control, which reads the trees entirely. Every probe counts as one query; transfer from a surrogate is covered by CAPGD and Online Resource~1 (Section~S11).

The transformed trace must keep the malicious semantics so that exploit delivery, credential cracking or service degradation still succeeds; we do not test that. Padding appends non-executable bytes (e.g.\ SSH ignore messages \cite{rfc4253}), and dilation lowers the burst intensity of volumetric floods. The NIDS is an ingress biflow sensor and the adversary controls only the client endpoint, so it changes the forward features $x_2,x_4,x_6,x_7,x_{11}$ and the derived rates, while the server-generated $x_3$ and $x_5$ are immutable. Chaff is sent as pure ACKs, which a TCP endpoint does not acknowledge (RFC 9293 \cite{rfc9293}), so $\Delta x_3=\Delta x_5=0$. \emph{Re-segmentation is different}: more forward data segments can elicit more server ACKs under delayed-ACK rules, which would raise $x_3$; the schema holds $x_3$ fixed, so the segmentation moves assume that the added segments do not change the backward direction. We flag this as a limitation of the schema-level feasibility notion.

The adversary minimizes the operational cost $\mathcal{C}(x^*,x_{\text{orig}})=\sum_k c_ks_k$, where $s_k$ is the amount of generator $k$ applied and $c_k$ its unit cost in latency, bandwidth and socket resources (Section~\ref{sec:cost}).

\section{The Flow-Field Evasion Framework}
\label{sec:framework}

\subsection{Geometry of the Flow Manifold and Its Generators}
\label{sec:geometry}

The physical flow set $\mathcal{M}_{\text{flow}}$ is a discrete lattice in $\mathbb{R}^{12}$. Its differential geometry is analyzed through the smooth relaxation $\widetilde{\mathcal{M}}_{\text{flow}}$ (Section~\ref{sec:formulation}), the zero set of the conservation residual
\begin{equation}
    \Psi(x)=\begin{bmatrix} x_8-10^6(x_4+x_5)/x_6\\ x_9-10^6(x_2+x_3)/x_6\\ x_{10}-x_4/x_2\end{bmatrix}=\mathbf{0}
\end{equation}
subject to the bounds listed there. At points where the Jacobian $J_\Psi$ has full rank 3, the tangent space is $\ker J_\Psi(x)$, so a physically realizable velocity satisfies the linearized conservation laws $\dot x_8=\tfrac{10^6(\dot x_4+\dot x_5)}{x_6}-\tfrac{x_8}{x_6}\dot x_6$, $\dot x_9=\tfrac{10^6(\dot x_2+\dot x_3)}{x_6}-\tfrac{x_9}{x_6}\dot x_6$ and $\dot x_{10}=\tfrac{\dot x_4-x_{10}\dot x_2}{x_2}$.

The attacker controls five elementary degrees of freedom, each generating a one-parameter flow whose push-forward is a tangent vector field ($\vec e_k$ is the $k$th basis vector of $\mathbb{R}^{12}$, indexed from 0):
\begin{equation}
\label{eq:generators}
\begin{aligned}
    \vec v_{\text{time}} &= \vec e_6 - \tfrac{x_8}{x_6}\vec e_8 - \tfrac{x_9}{x_6}\vec e_9, \qquad
    \vec v_{\text{pad}} = \vec e_4 + \tfrac{10^6}{x_6}\vec e_8 + \tfrac{1}{x_2}\vec e_{10},\\
    \vec v_{\text{chaff}} &= \vec e_2 + L_0\vec e_4 + \tfrac{L_0 10^6}{x_6}\vec e_8 + \tfrac{10^6}{x_6}\vec e_9 + \tfrac{L_0-x_{10}}{x_2}\vec e_{10} + H_0\vec e_{11},\\
    \vec v_{\text{seg}} &= \vec e_2 + \tfrac{10^6}{x_6}\vec e_9 - \tfrac{x_{10}}{x_2}\vec e_{10} + H_0\vec e_{11},\qquad
    \vec v_{\text{win}} = \mathbb{I}_{x_1=6}\,\operatorname{sat}\!\left(\tfrac{W_{\text{target}}-x_7}{\epsilon_{\text{win}}}\right)\vec e_7.
\end{aligned}
\end{equation}
Temporal dilation adds delay without bytes or packets, padding adds forward dummy bytes, chaff adds forward header-only probes, booked as $L_0$ bytes of the length column, and window morphing moves the client window toward a target $W_{\text{target}}$. The probe sizes are $L_0=40$, $H_0=20$ for empty TCP segments \cite{rfc9293}, $L_0=52$, $H_0=32$ with TCP timestamps \cite{rfc7323}, and $L_0=28$, $H_0=8$ for UDP \cite{rfc768}. The new generator, \emph{re-segmentation} (an MSS clamp), carries the same forward payload in more, smaller packets: the packet count and the header total grow by $H_0$ per packet, while the byte total and the duration are unchanged and the mean length falls. It is the chaff field with $L_0=0$, hence $\vec v_{\text{seg}}=\vec v_{\text{chaff}}-L_0\vec v_{\text{pad}}$ component-wise: the five fields commute but are \emph{not} linearly independent, and re-segmentation lies in the span of chaff and padding but outside their nonnegative cone (it corresponds to chaff with negative padding). An integer move of $K$ packets is exactly the time-$K$ flow of this field ($x_{10}=x_4/(x_2+K)$). A sixth generator, TCP option padding ($\vec v_{\text{hdr}}=\vec e_{11}$), is implemented but is not part of the frozen configuration. The window field is globally Lipschitz; the implementation uses its limit $\epsilon_{\text{win}}\to0$, which lands exactly on $W_{\text{target}}$.

\subsection{Commutativity, Invariants and the Lie Cone}
\label{sec:invariants}

\begin{proposition}[Commutativity of the physical generators]
\label{prop:lie_commutativity}
Because the protocol $x_1\in\{6,17\}$ is immutable, $\widetilde{\mathcal{M}}_{\text{flow}}$ decomposes into a TCP and a UDP component on which $\mathbb{I}_{x_1=6}$, $L_0$ and $H_0$ are constant along trajectories (FFE-v2 fixes $H_0=32$ iff the mean forward header $x_{11}/x_2\ge26$ \emph{of the original flow}; FFE-v1 re-evaluates it on the current flow, where injecting packets of header size $H_0$ moves that mean toward $H_0$ and never across the threshold). On each component the fields $\vec v_{\text{time}},\vec v_{\text{pad}},\vec v_{\text{chaff}},\vec v_{\text{seg}}$ are $C^\infty$, $\vec v_{\text{win}}$ is Lipschitz, and
\begin{equation}
    [\vec v_i,\vec v_j]\equiv\mathbf{0}\qquad\text{for all }i,j\in\{\text{time},\text{pad},\text{chaff},\text{seg},\text{win}\},
\end{equation}
so the fields span an Abelian Lie algebra, of dimension 4 on TCP (3 on UDP) because $\vec v_{\text{seg}}$ is a linear combination of $\vec v_{\text{chaff}}$ and $\vec v_{\text{pad}}$. The same holds with $\vec v_{\text{hdr}}$.
\end{proposition}
The proof is an entry-wise check of the four new brackets involving $\vec v_{\text{seg}}$ on top of the six of FFE-v1 (Online Resource~1, Section~S1, Lemma S1); the window field acts only on $x_7$, which no other field reads or writes. The admissible cone is $\mathcal{K}_{\text{net}}(x)=\{\sum_k u_k\vec v_k(x)\mid u_k\ge0\}\subset T_x\widetilde{\mathcal{M}}_{\text{flow}}$.

\begin{theorem}[Invariant preservation]
\label{thm:invariant_universal}
Let $x_0\in\widetilde{\mathcal{M}}_{\text{flow}}$ and $u\in L^\infty_{\text{loc}}(\mathbb{R}_{\ge0},\mathbb{R}^5_{\ge0})$ satisfy the admissibility rules: $u_{\text{pad}}\le(1500-L_0)u_{\text{chaff}}+1500u_{\text{seg}}$ whenever $x_4=1500x_2$; $u_{\text{time}}=0$ whenever $x_6=1.1\times10^8$; and $u_{\text{seg}}=0$ whenever $x_{10}\le1$ (a re-segmented packet carries at least one payload byte). Then the solution of $\dot x=\sum_ku_k\vec v_k(x)$, $x(0)=x_0$, exists, is unique, is global and stays in $\widetilde{\mathcal{M}}_{\text{flow}}$.
\end{theorem}
\noindent\textit{Proof sketch.} Differentiating $\Psi$ along the flow gives $\dot\Psi=A(t)\Psi$ with $A=\operatorname{diag}\!\big(-\tfrac{u_{\text{time}}}{x_6},-\tfrac{u_{\text{time}}}{x_6},-\tfrac{u_{\text{chaff}}+u_{\text{seg}}}{x_2}\big)$ locally integrable, so $\Psi\equiv0$ by Carath\'eodory's theorem \cite{caratheodory1918vorlesungen}. All remaining bounds are affine and satisfy the Nagumo--Brezis condition \cite{nagumo1942lage, brezis1970characterization} on every face under the admissibility rules, which gives forward invariance, existence and uniqueness. The full proof, including the new case of $\vec v_{\text{seg}}$ at the MTU face and its effect on the header-total bound, is in Online Resource~1, Section~S1. Under discrete stepping the closed-form operator $\mathcal{P}_{\mathcal{M}}$ is a retraction (exact on feasible inputs, not a metric projection) that removes truncation drift and enforces integrality.

The theorem concerns the continuous relaxation and the schema, not side effects at a real endpoint (Section~\ref{sec:formulation}).

\subsection{Operational Cost}
\label{sec:cost}

The unit costs are $c_1=10^{-4}$ per second of induced delay, $c_2=10^{-2}$ per byte of padding, $c_3=1.0$ per added packet and $c_4=10^{-3}$ per normalized window displacement $|W_{\text{target}}-x_7|/65535$. A re-segmented packet is priced exactly like a chaff packet, and header-option bytes are not priced. Because the generators commute, the displacement on the interior of the manifold depends only on the integrated amounts $s_k=\int u_k\,dt\ge0$, and the cost $\mathcal{C}(s)=\sum_kc_ks_k$ is linear and path-independent there (not at the boundary faces). The cost model is the one FFE steers by. The boundary $F(x)=0$ of a tree ensemble is piecewise constant and non-convex, so the evasion problem, minimizing $\mathcal{C}$ subject to $F(x(T^*))<0$ and the dynamics above, is solved heuristically (Section~\ref{sec:algorithm}).

\subsection{Algorithmic Realization: FFE-v2}
\label{sec:algorithm}

FFE-v1 took one fixed-size probe per generator and iteration (Online Resource~1, Algorithm~S1). Validation diagnostics (Online Resource~1, Section~S3) showed where that stalls: on UNSW-NB15 the score plateaus until splits on the packet count, mean length and packet rate are crossed jointly, which needs payload-free packets in magnitudes of tens to hundreds, and elsewhere the single window target lay on the wrong side of the model's borders. FFE-v2 changes five things (Online Resource~1, Algorithm~S5); all other machinery, in particular the projection $\mathcal{P}_{\mathcal{M}}$ that follows every move (Online Resource~1, Algorithm~S2), is unchanged.



\paragraph{Probing, targets and selection.} Every generator is probed on a geometric ladder of integer magnitudes $a_k=\min(a_0r^k,\,a_{\max})$ spanning its whole feasible range (delay from 1~ms to the 110~s cap, padding from 16~B to the MTU headroom, one to 4,096 chaff or segmented packets), so a plateau of the piecewise-constant score is crossed within one iteration; identical projected candidates are probed once. The window generator tries the six most frequent benign TCP windows of the attacker's benign sample (the training split), in both directions, instead of 64,240~B; a variant with generic operating-system defaults needs no data. Progress is the decrease of the log-odds of the returned score, which does not saturate where float32 probabilities near 1 tie. If any probe crosses 0.5, the cheapest crossing under the operational cost of Section~\ref{sec:cost} ends the search and four bisection probes shrink its magnitude (each probe is a scored flow, so success is never assumed); otherwise the move with the largest decrease per unit of marginal cost is taken, switching to pure score minimization after three iterations or when no move is cost-efficient.

\paragraph{Optional snapping (Tier 1).} When no probe helps, threshold crossings of \emph{all} mutable features (eight per feature and direction) are converted to closed-form generator magnitudes, including exact inversions of the rate and mean thresholds. This is the only use of the trees' thresholds; the headline variant does not snap.

\paragraph{Queries and configuration.} Every scored flow, including probes scored in a batch, is one query; observed means are 28--38 on every cell except CIC-IDS2017 CatBoost (109), where the budget is exhausted on unreachable flows. The hyper-parameters were chosen on a validation split over 58 configurations (Online Resource~1, Table~S1) and frozen before the single test evaluation: ladder ratio, refinement probes and window pool size are in \texttt{results/ffev2\_config.json}, and the iteration budget $M=8$, the relax rule $R=3$ and the ladder range are the defaults of \texttt{src/attacks/ffe\_v2.py}. Beam search and option padding gave no gain worth their queries.

\section{Experimental Protocol}
\label{sec:experiments}
\label{sec:empirical_evaluation}

\subsection{Corpora, Targets and Attackers}
\label{sec:setup}

Single-testbed evaluations have limited external validity~\cite{cantone2024crossdataset}, and attack evaluations that ignore realism overstate the risk~\cite{apruzzese2022modeling, apruzzese2023real}, so we use three widely studied corpora in full: CSE-CIC-IDS2018~\cite{sharafaldin2018toward} (nine daily files, seven capped at 1,048,575 rows in the distribution), CIC-IDS2017~\cite{sharafaldin2018toward} (five daily files) and UNSW-NB15~\cite{moustafa2015unsw} (the official training and testing partition, 257,673 of the 2.54 million records). Every record is mapped onto the 12-feature schema of Section~\ref{sec:formulation}, invalid records are dropped, and every flow is projected onto the manifold \emph{before training and testing} (Online Resource~1, Section~S12); this pre-projection clamps byte totals above 1500 times the packet count, durations above 110~s and UDP header sizes, so the learned models never see such values. Records with identical features and label are then collapsed so that no flow appears on both sides of the split~\cite{engelen2021troubleshooting, lanvin2023errors}, which removes 39.1\%, 20.5\% and 42.7\% of the valid records and leaves 4,968,389, 1,658,869 and 123,508 flows (Online Resource~1, Table~S20). Records with identical features but different labels (149,526, 22 and 1,040) are kept as real ambiguity of the schema. The split into training, validation and test partitions is 70\%/10\%/20\%, stratified by label and family, random, and not grouped by time or host, so near-duplicates may fall on both sides and accuracies are optimistic.

The targets are XGBoost (CUDA), LightGBM (each 60 trees of depth 6, learning rate 0.1) and CatBoost (60 oblivious trees of depth 6). \emph{FFE-v2} is score-only; its snapping variant also reads thresholds. \emph{FFE-v1} (Online Resource~1, Algorithm~S1) is the reference baseline and reads thresholds for its snapping step. The \emph{projected white-box control} reads all trees: each step moves one feature across one split threshold, choosing the move that lowers the exact ensemble margin most, with every candidate projected onto the manifold before scoring; it counts two queries (the tree reads are free) and is cost-blind. For XGBoost it is a greedy leaf-tuple-style attack in the spirit of Zhang et al.~\cite{zhang2020efficient}; for LightGBM and CatBoost it is our generic version on the raw ensemble margin. It is not LT-Attack (whose official code needs a Bazel/C++ build we did not set up) and we did not run Cascioli et al.~\cite{cascioli2024faster}. On XGBoost we also compare with seven earlier baselines: a naive coordinate descent, \emph{ZOO-NIDS}~\cite{chen2017zoo}, \emph{HopSkipJump-NIDS}~\cite{chen2020hopskipjump}, \emph{MOEVA}~\cite{simonetto2024constrained, ghamizi2020moeva}, \emph{PSO}~\cite{han2021traffic} with our projection, and our re-implementations of \emph{CAPGD}~\cite{simonetto2024constrained} and a \emph{NIDSGAN-style generator}~\cite{zolbayar2022nidsgan}, which fit a small MLP surrogate on the training flows (surrogate computations are not counted). CAPGD was run on XGBoost in the single-split test and, in the seed study, also on LightGBM and CatBoost (three seeds); baselines use default budgets.

\subsection{Protocol, Metrics and Cohorts}
\label{sec:protocol}

All design choices and hyper-parameters of FFE-v2 were tuned (tuning queries are not counted) on cohorts of up to 1,000 detected malicious flows from the 10\% \emph{validation} split; the configuration was frozen and evaluated once on the benchmark's 1,000-flow \emph{test} cohorts (the identical flows used for FFE-v1 and the baselines, verified by the runner). This is not a clean hold-out: the FFE-v1 failures that motivated FFE-v2 were first observed on these test cohorts, re-segmentation was motivated by the control's moves on the same corpora, and the window pool uses the attacker's benign training sample. The five-seed study (retrained models, fresh cohorts, no tuning) is the cleaner confirmation. The cohorts are drawn once, with a fixed seed, from the malicious test flows that both XGBoost and LightGBM detect when clean (77.6\%, 99.8\% and 87.0\% of the test attacks of CSE-CIC-IDS2018, CIC-IDS2017 and UNSW-NB15), so all results are conditional on detection. For CatBoost the cohort is restricted to the flows CatBoost itself detects ($N=1000$, 997 and 983), so every success is a genuine evasion of that model. MOEVA and PSO need 183--445 queries per flow and attack the first 200 flows. FFE-v1 on XGBoost and LightGBM and the baselines are the saved benchmark outcomes; FFE-v1 on CatBoost and the control on LightGBM and CatBoost were run for this study.

ASR is the fraction of attacks with $P(\text{Malicious}\mid x^*)<0.5$; PFR the fraction of evasive candidates that satisfy the schema invariants; Real ASR counts an attack only when it both evades and is feasible. Feasibility is the schema verifier of Section~\ref{sec:formulation} (immutability, monotonicity, MTU and header bounds, 5\% tolerance on the coupled rates, identical for every method), which holds for FFE by construction; it does not check the wire-byte gap $x_4\ge x_{11}+20x_2$, which re-segmentation would violate if $x_4$ counted headers. ``100\% schema-feasible'' is relative to this verifier. We report Wilson 95\% intervals and exact McNemar tests on paired outcomes; no multiplicity correction is applied except where stated (nine cells, so individual $p$-values should be read with that in mind), and pooled seed tests treat flows as independent. Cost is the median over a method's own feasible evasions of the operational cost of Section~\ref{sec:cost} and of the normalized distance $D_{L_1}$; these medians cover different flows, so we add paired medians on flows two methods both evade. For the multi-seed study every seed retrains the targets on a new stratified split and redraws a cohort of up to 300 flows detected by XGBoost (LightGBM and CatBoost are attacked on the part of that cohort they also detect); the frozen configuration is used without tuning, with five seeds for XGBoost and three for LightGBM and CatBoost.

\section{Attack Results}
\label{sec:results}

\subsection{Benchmark Against XGBoost}
\label{sec:master}

\begin{table}[!htbp]
\centering
\caption{Evasion benchmark against CUDA XGBoost on the full corpora. All methods except MOEVA and PSO attack the same 1,000 detected malicious test flows; MOEVA and PSO attack the first 200 of them. PFR: share of evasive candidates that are schema-feasible; Op.\ cost and $D_{L_1}$: medians over the method's own feasible evasions. FFE-v1 reads thresholds, the control reads the trees, FFE-v2 (bold) reads scores only.}\label{tab:master_evasion_benchmark}
\resizebox{\textwidth}{!}{\begin{varwidth}{\maxdimen}\noindent%
\begin{tabular}{@{}lccccccc@{}}
\toprule
\textbf{Method} & \textbf{ASR} & \textbf{PFR} & \textbf{Real ASR [95\% CI]} & \textbf{Queries} & \textbf{Op.\ cost} & \textbf{$D_{L_1}$} & \textbf{ms/flow} \\
\midrule
\multicolumn{8}{@{}l}{\textbf{CSE-CIC-IDS2018} (XGBoost accuracy 97.30\%, test attacks 118,963, detected by both models 92,275)} \\
\midrule
Unconstrained feature space & 36.8 & 88.0 & 32.4 [29.6, 35.4] & 66.8 & 0.000 & 0.000 & 48.0 \\
ZOO-NIDS \cite{chen2017zoo} & 30.3 & 83.2 & 25.2 [22.6, 28.0] & 116.7 & 0.000 & 0.000 & 91.0 \\
HopSkipJump-NIDS \cite{chen2020hopskipjump} & 100.0 & 0.0 & 0.0 [0.0, 0.4] & 47.0 & -- & -- & 30.3 \\
MOEVA (search stage of CAA) \cite{simonetto2024constrained, ghamizi2020moeva} & 40.5 & 70.4 & 28.5 [22.7, 35.1] & 330.1 & 0.092 & 0.001 & 253.2 \\
PSO (Traffic Manipulator, projected) \cite{han2021traffic} & 54.0 & 100.0 & 54.0 [47.1, 60.8] & 182.9 & 16.3 & 0.061 & 142.0 \\
CAPGD (surrogate) \cite{simonetto2024constrained} & 88.5 & 100.0 & 88.5 [86.4, 90.3] & 3.4 & 46.1 & 0.824 & 15.4 \\
NIDSGAN-style \cite{zolbayar2022nidsgan} & 61.3 & 100.0 & 61.3 [58.2, 64.3] & 2.5 & 48.0 & 0.836 & 7.1 \\
FFE-v1 (reads thresholds) & 99.5 & 100.0 & 99.5 [98.8, 99.8] & 16.5 & 0.001 & 0.116 & 15.8 \\
White-box control (projected, reads trees) & 99.5 & 100.0 & 99.5 [98.8, 99.8] & 2.0 & 0.000 & 0.120 & 19.1 \\
FFE-v2 + snapping (reads thresholds) & 99.5 & 100.0 & 99.5 [98.8, 99.8] & 30.8 & 0.000 & 0.006 & 36.0 \\
\textbf{FFE-v2 (score-only)} & 99.5 & 100.0 & \textbf{99.5 [98.8, 99.8]} & 30.3 & 0.000 & 0.006 & 44.6 \\
\midrule
\multicolumn{8}{@{}l}{\textbf{CIC-IDS2017} (XGBoost accuracy 99.96\%, test attacks 79,281, detected by both models 79,150)} \\
\midrule
Unconstrained feature space & 64.8 & 34.4 & 22.3 [19.8, 25.0] & 37.6 & 0.000 & 0.000 & 30.6 \\
ZOO-NIDS \cite{chen2017zoo} & 62.2 & 35.0 & 21.8 [19.4, 24.5] & 79.9 & 0.000 & 0.000 & 58.5 \\
HopSkipJump-NIDS \cite{chen2020hopskipjump} & 100.0 & 0.0 & 0.0 [0.0, 0.4] & 47.0 & -- & -- & 36.5 \\
MOEVA (search stage of CAA) \cite{simonetto2024constrained, ghamizi2020moeva} & 64.0 & 59.4 & 38.0 [31.6, 44.9] & 225.9 & 0.059 & 0.005 & 201.0 \\
PSO (Traffic Manipulator, projected) \cite{han2021traffic} & 35.0 & 100.0 & 35.0 [28.7, 41.8] & 251.4 & 15.2 & 0.137 & 203.3 \\
CAPGD (surrogate) \cite{simonetto2024constrained} & 89.5 & 100.0 & 89.5 [87.4, 91.3] & 3.2 & 15.0 & 0.305 & 15.1 \\
NIDSGAN-style \cite{zolbayar2022nidsgan} & 65.3 & 100.0 & 65.3 [62.3, 68.2] & 2.4 & 24.6 & 0.490 & 5.8 \\
FFE-v1 (reads thresholds) & 99.2 & 100.0 & 99.2 [98.4, 99.6] & 14.6 & 0.001 & 0.289 & 10.7 \\
White-box control (projected, reads trees) & 100.0 & 100.0 & 100.0 [99.6, 100.0] & 2.0 & 0.000 & 0.045 & 11.9 \\
FFE-v2 + snapping (reads thresholds) & 100.0 & 100.0 & 100.0 [99.6, 100.0] & 28.3 & 0.001 & 0.187 & 26.5 \\
\textbf{FFE-v2 (score-only)} & 100.0 & 100.0 & \textbf{100.0 [99.6, 100.0]} & 28.3 & 0.001 & 0.187 & 37.1 \\
\midrule
\multicolumn{8}{@{}l}{\textbf{UNSW-NB15} (XGBoost accuracy 88.51\%, test attacks 11,113, detected by both models 9,673)} \\
\midrule
Unconstrained feature space & 21.6 & 0.9 & 0.2 [0.1, 0.7] & 48.5 & 0.082 & 0.000 & 49.1 \\
ZOO-NIDS \cite{chen2017zoo} & 20.6 & 0.0 & 0.0 [0.0, 0.4] & 130.1 & -- & -- & 96.5 \\
HopSkipJump-NIDS \cite{chen2020hopskipjump} & 8.1 & 0.0 & 0.0 [0.0, 0.4] & 47.0 & -- & -- & 40.6 \\
MOEVA (search stage of CAA) \cite{simonetto2024constrained, ghamizi2020moeva} & 13.0 & 100.0 & 13.0 [9.0, 18.4] & 445.0 & 0.092 & 0.002 & 339.2 \\
PSO (Traffic Manipulator, projected) \cite{han2021traffic} & 42.0 & 100.0 & 42.0 [35.4, 48.9] & 228.4 & 21.6 & 0.060 & 164.3 \\
CAPGD (surrogate) \cite{simonetto2024constrained} & 21.1 & 100.0 & 21.1 [18.7, 23.7] & 16.2 & 53.2 & 1.158 & 73.4 \\
NIDSGAN-style \cite{zolbayar2022nidsgan} & 11.3 & 100.0 & 11.3 [9.5, 13.4] & 4.5 & 34.0 & 1.211 & 11.2 \\
FFE-v1 (reads thresholds) & 57.6 & 100.0 & 57.6 [54.5, 60.6] & 93.7 & 5.610 & 0.027 & 89.1 \\
White-box control (projected, reads trees) & 90.2 & 100.0 & 90.2 [88.2, 91.9] & 2.0 & 213.0 & 0.170 & 26.4 \\
FFE-v2 + snapping (reads thresholds) & 98.5 & 100.0 & 98.5 [97.5, 99.1] & 37.9 & 36.0 & 0.121 & 81.8 \\
\textbf{FFE-v2 (score-only)} & 98.3 & 100.0 & \textbf{98.3 [97.3, 98.9]} & 36.2 & 40.0 & 0.121 & 68.0 \\
\bottomrule
\end{tabular}
%
\end{varwidth}}
\end{table}

Table~\ref{tab:master_evasion_benchmark} compares FFE-v2 with FFE-v1, the projected white-box control and the seven earlier baselines. Feature-space success is not feasibility: HopSkipJump-NIDS reaches 100.0\% ASR on two corpora with no feasible evasive flow, and the naive attack and ZOO-NIDS keep 83--88\% of their evasions on CSE-CIC-IDS2018 at zero cost, consistent with moving only coupled rates within the 5\% tolerance, so that check flatters them.

FFE-v2 reaches a Real ASR of 99.5\%, 100.0\% and 98.3\% on CSE-CIC-IDS2018, CIC-IDS2017 and UNSW-NB15. It is significantly above all seven earlier baselines on every corpus (21 exact McNemar tests, largest Holm-adjusted $p=8\times10^{-28}$); the closest competitor, CAPGD (88.5\%, 89.5\%, 21.1\%), is beaten on 110, 105 and 772 flows against none. Against the stronger references the picture is more precise. On the two easy corpora our greedy white-box control, which reads the trees, ties FFE-v2 (99.5\% and 100.0\%, identical flows); the five CSE-CIC-IDS2018 flows that FFE-v2, FFE-v1 and the control all miss are DDoS flows that no tested move reached. On UNSW-NB15 FFE-v2 exceeds the control (98.3\% against 90.2\%; 89 flows against 8, $p=2\times10^{-18}$) and FFE-v1 (57.6\%; 412 against 5). FFE-v1 had been 32.6 points \emph{below} the control there, and the difference is the repair described in Section~\ref{sec:algorithm}: on the validation flows that FFE-v1 missed, the control's successful moves raise the packet count by a median of 213 and in 82\% of the cases leave the byte total unchanged, which is the re-segmentation move that FFE-v1 could not express (Online Resource~1, Section~S3). FFE-v2 with snapping, which reads thresholds, is statistically indistinguishable on XGBoost (98.5\% on UNSW-NB15; $p=0.5$), so the score-only variant is the headline.

The methods pay differently. CAPGD and the NIDSGAN-style generator need far fewer target queries (3.2--16.2 and 2.4--4.5) because their gradient work runs on an uncounted surrogate, and the control needs two. FFE-v2 averages 30.3, 28.3 and 36.2 queries, about twice FFE-v1 on the easy corpora (16.5 and 14.6) and well below it on UNSW-NB15 (93.7). On the same 200 flows it needs 6 to 12 times fewer queries than PSO and MOEVA, although query counts are not strictly like-for-like across score, decision and surrogate methods. The median operational cost of FFE-v2's feasible evasions is $3.0\times10^{-5}$, $5.5\times10^{-4}$ and 40.0 units, against $5.7\times10^{-4}$, $8.6\times10^{-4}$ and 5.6 for FFE-v1; on UNSW-NB15 this partly reflects who succeeds (FFE-v1 evades only the cheap flows; Section~\ref{sec:ablation}), and the median of 40 is dominated by tens to hundreds of chaff or re-segmented packets priced at one unit each. These comparisons are specific to the cost model, which FFE steers by.

\subsection{Three Libraries and Three Corpora}
\label{sec:multidataset_results}

\begin{table}[!htbp]
\centering
\caption{Real ASR (\%) on the 1,000-flow test cohorts (CatBoost: flows it detects), with clean test accuracy and mean queries per flow ($q$). Single frozen evaluation.}\label{tab:crosslib}
\resizebox{\textwidth}{!}{\begin{varwidth}{\maxdimen}\noindent%
\begin{tabular}{@{}llccccccrr@{}}
\toprule
\textbf{Corpus} & \textbf{Target} & \textbf{$N$} & \textbf{Acc.} & \textbf{FFE-v1} & \textbf{Control} & \textbf{FFE-v2} & \textbf{+snap} & \textbf{$q$ v1} & \textbf{$q$ v2} \\
\midrule
CSE-CIC-IDS2018 & XGBoost & 1000 & 97.30 & 99.5 & 99.5 & \textbf{99.5} & 99.5 & 16.5 & 30.3 \\
 & LightGBM & 1000 & 97.28 & 99.5 & 99.5 & \textbf{99.5} & 99.5 & 20.2 & 31.1 \\
 & CatBoost & 1000 & 97.26 & 59.3 & 99.5 & \textbf{99.5} & 99.5 & 52.5 & 29.1 \\
\midrule
CIC-IDS2017 & XGBoost & 1000 & 99.96 & 99.2 & 100.0 & \textbf{100.0} & 100.0 & 14.6 & 28.3 \\
 & LightGBM & 1000 & 99.96 & 35.1 & 36.6 & \textbf{96.0} & 96.0 & 120.1 & 32.4 \\
 & CatBoost & 997 & 99.85 & 34.8 & 36.4 & \textbf{36.4} & 36.4 & 76.3 & 109.0 \\
\midrule
UNSW-NB15 & XGBoost & 1000 & 88.51 & 57.6 & 90.2 & \textbf{98.3} & 98.5 & 93.7 & 36.2 \\
 & LightGBM & 1000 & 88.47 & 62.4 & 88.1 & \textbf{96.7} & 99.2 & 89.0 & 38.3 \\
 & CatBoost & 983 & 86.67 & 70.1 & 56.8 & \textbf{97.5} & 97.7 & 67.9 & 38.0 \\
\bottomrule
\end{tabular}
%
\end{varwidth}}
\end{table}

Table~\ref{tab:crosslib} extends the comparison to LightGBM and CatBoost. FFE-v2 is at least as successful as our greedy projected control and as FFE-v1 in all nine cells. It is significantly higher than the control in four (UNSW-NB15 for all three libraries, $p\le5\times10^{-15}$, and CIC-IDS2017 LightGBM, 594 flows against none) and than FFE-v1 in seven (exact McNemar; the weakest is CIC-IDS2017 XGBoost, 8 flows against none, $p=0.008$); the remaining cells tie with identical flows. The two failure modes of FFE-v1 are repaired. On CSE-CIC-IDS2018 CatBoost FFE-v1 evaded 59.3\% because its single window target (64,240~B) lies on the wrong side of CatBoost's borders, whereas on the validation cohort the control evades 411 of FFE-v1's 414 failures with a window move alone, and FFE-v2 reaches 99.5\%. On CIC-IDS2017 LightGBM FFE-v1 reached 35.1\% for the same reason: 64,240~B almost never occurs in the CIC-IDS2017 training flows (0.08\% of benign and 0.015\% of attack TCP flows), so each ensemble places its splits arbitrarily in that near-empty interval (XGBoost routes it to the benign side, LightGBM keeps it malicious; Online Resource~1, Section~S9). Drawing the targets from the attacker's benign sample lifts LightGBM to 96.0\%, and the residual is concentrated in DoS flows (90.5\% evaded), which we did not examine further.

\paragraph{A ceiling, not a win.} CatBoost on CIC-IDS2017 stays at 36.4\% for FFE-v2 and the control (34.8\% for FFE-v1). The decision for DDoS and DoS flows rests on attacker-immutable features (backward packets, backward bytes and destination port carry 77\% of CatBoost's gain), and a 20,000-sample multi-start search over all generators found no crossing for any of 37 sampled DDoS and 40 sampled DoS validation flows (3 of 80 sampled failed flows were reachable, all port scans; \texttt{results/ffev2\_report.md}). FFE-v2's own Real ASR on these families is 0.0\% (DDoS) and 5.0\% (DoS). We found no crossing in a multi-start search, which cannot prove unreachability; the number describes the model's reliance on immutable fields, not attack strength, and the unreachable flows make FFE-v2 spend its whole budget (109 queries).

\paragraph{Snapping and per-family results.} Threshold snapping adds up to 2.5 points on UNSW-NB15 (LightGBM, 96.7\% to 99.2\%, $p=6\times10^{-8}$) and at most 0.2 point elsewhere. Per-family rates are in Online Resource~1, Table~S5; on UNSW-NB15 every XGBoost family with more than 20 flows is above 90\%.

\subsection{Retraining Variance}
\label{sec:seeds}

\begin{table}[!htbp]
\caption{Real ASR (\%, mean $\pm$ standard deviation over seeds) of the frozen attacks; each seed retrains the target and draws a new cohort of up to 300 detected flows. Five seeds for XGBoost, three for LightGBM and CatBoost. $q$: mean queries of FFE-v2.}\label{tab:seeds}
\centering
\resizebox{\textwidth}{!}{\begin{varwidth}{\maxdimen}\noindent%
\begin{tabular}{@{}llcccccr@{}}
\toprule
\textbf{Corpus} & \textbf{Target} & \textbf{Seeds} & \textbf{FFE-v2} & \textbf{FFE-v1} & \textbf{Control} & \textbf{CAPGD} & \textbf{$q$ (v2)} \\
\midrule
CSE-CIC-IDS2018 & XGBoost & 5 & \textbf{99.7 $\pm$ 0.3} & 99.7 $\pm$ 0.3 & 99.7 $\pm$ 0.3 & 97.7 $\pm$ 2.4 & 31.1 \\
 & LightGBM & 3 & \textbf{99.7 $\pm$ 0.3} & 99.7 $\pm$ 0.3 & 99.7 $\pm$ 0.3 & 98.9 $\pm$ 1.1 & 31.2 \\
 & CatBoost & 3 & \textbf{99.7 $\pm$ 0.3} & 56.4 $\pm$ 4.3 & 99.7 $\pm$ 0.3 & 94.4 $\pm$ 2.5 & 29.2 \\
\midrule
CIC-IDS2017 & XGBoost & 5 & \textbf{100.0 $\pm$ 0.0} & 97.3 $\pm$ 1.5 & 100.0 $\pm$ 0.0 & 82.4 $\pm$ 11.9 & 29.1 \\
 & LightGBM & 3 & \textbf{95.6 $\pm$ 0.2} & 36.1 $\pm$ 4.2 & 77.4 $\pm$ 31.6 & 50.0 $\pm$ 2.4 & 30.3 \\
 & CatBoost & 3 & \textbf{36.7 $\pm$ 3.5} & 35.1 $\pm$ 3.1 & 36.5 $\pm$ 3.3 & 23.0 $\pm$ 2.8 & 122.1 \\
\midrule
UNSW-NB15 & XGBoost & 5 & \textbf{97.7 $\pm$ 2.5} & 58.3 $\pm$ 4.3 & 86.0 $\pm$ 8.0 & 31.5 $\pm$ 2.8 & 38.4 \\
 & LightGBM & 3 & \textbf{98.9 $\pm$ 1.0} & 57.2 $\pm$ 5.5 & 90.9 $\pm$ 8.2 & 30.4 $\pm$ 4.9 & 37.0 \\
 & CatBoost & 3 & \textbf{99.3 $\pm$ 0.3} & 72.1 $\pm$ 0.9 & 74.4 $\pm$ 5.6 & 30.1 $\pm$ 6.2 & 37.2 \\
\bottomrule
\end{tabular}
%
\end{varwidth}}
\end{table}

A single split and cohort cannot show how much of the ranking is luck. Table~\ref{tab:seeds} repeats the comparison without any tuning. FFE-v2 is never below FFE-v1, the control or CAPGD in any seed of any cell, and strictly above them in most. On UNSW-NB15 XGBoost its mean is 97.7\% against 58.3\% (FFE-v1), 86.0\% (control) and 31.5\% (CAPGD); pooled over the five seeds it evades 188 flows that the control misses and misses 12 that the control evades ($p=8\times10^{-42}$). The control, our own implementation, is erratic where the model depends on a single move: on CIC-IDS2017 LightGBM its Real ASR is 95.7\%, 41.0\% and 95.7\% over the three seeds (mean 77.4\%), whereas FFE-v2 holds 95.6\% $\pm$ 0.2. CatBoost on CIC-IDS2017 remains at the ceiling (36.7\% $\pm$ 3.5). Pooled McNemar tests and win--tie--loss counts are in Online Resource~1, Table~S6. Three seeds for LightGBM and CatBoost are few, pooled $p$-values treat flows as independent, and nine cells are tested without multiplicity correction. The seeds vary the split and the cohort but the configuration was tuned on the main split's validation cohorts; they are the cleaner confirmation, not an independent tuning test.

\section{What Drives the Attack: Ablations, Cost and Packet-Level Replay}
\label{sec:ablation}

\paragraph{Ablations on the validation split.} Online Resource~1 (Table~S1) varies one factor at a time over nine cells. The data-driven window pool is decisive on the CSE-CIC-IDS2018 and CIC-IDS2017 corpora (removing it drops mean Real ASR from 91.6\% to 60.6\%), chaff matters on UNSW-NB15 for LightGBM and CatBoost (59.9\% and 63.9\% without it), and the ladder must span the range. Re-segmentation adds little on validation once chaff and padding exist (96.6\%, 96.9\% and 97.2\% without it against 96.7\%, 97.5\% and 97.4\% on UNSW-NB15), yet 39\%, 16\% and 24\% of FFE-v2's successful UNSW-NB15 flows use it (Online Resource~1, Table~S4).

\paragraph{Cost.} FFE-v2 is not uniformly cheaper than FFE-v1: on UNSW-NB15 its paired median operational cost on the flows both evade is higher (6.7 against 5.8, 14.6 against 7.5 and 29.1 against 9.8 for XGBoost, LightGBM and CatBoost), and against the control it is cheaper for XGBoost and LightGBM (32.5 against 213 and 35.8 against 49) but not for CatBoost (22.1 against 10.5; Online Resource~1, Tables~S2 and~S3).

\paragraph{Packet-level replay.} To check that the generators correspond to realizable packet operations we expanded the action ledger of each successful flow of the first 300 test flows per corpus (XGBoost) into a synthetic packet trace, re-extracted the twelve features with a CICFlowMeter-style extractor (our re-implementation) and re-scored the target (Online Resource~1, Section~S8, Table~S15). When chaff packets carry the booked bytes, or empty chaff is compensated by padding, 100\% of the replayed flows reproduce the schema vector within 1\%, stay feasible and stay evasive, including the 118 UNSW-NB15 flows that use re-segmentation (replayed ASR 99.3\%, 100.0\% and 98.3\%). If chaff packets are empty and the bytes the schema credits to them are not sent, the easy corpora are unchanged but UNSW-NB15 retains only 69.2\% of its evasions (replayed ASR 68.0\%), the same sensitivity as for FFE-v1. The replay validates the action semantics and the schema's internal consistency only: the traces are synthetic, no TCP state machine or server reaction is modelled (in particular whether more forward segments elicit more server ACKs), and nothing was sent over a network.

\section{A Monotone-Gate Defense with a Conditional Certificate}
\label{sec:defense}

The attacks suggest what a defender can and cannot promise. The schema lets the attacker only \emph{increase} forward packets, forward bytes, duration and header bytes, hold the port, protocol and backward direction fixed, and move the window, the two rates and the mean length freely or as functions of the others. A score that never decreases when those four columns increase, that does not read the free or derived columns, and whose other inputs stay fixed, cannot be lowered by a feasible move. This section builds a defense on that observation, evaluates it, and states where it stops. The monotone-model idea is due to Incer et al.~\cite{incer2018adversarially}; what is new here is its use on coupled flow records, where the derived columns must be hidden.

\subsection{Design}
\label{sec:defense_design}

\paragraph{Component 1: monotone gate.} A tree ensemble that (i) is never shown the window, the two rates and the mean length ($x_7,x_8,x_9,x_{10}$), and (ii) carries monotone constraints (non-decreasing malicious score) on the four increase-only columns $x_2,x_4,x_6,x_{11}$, with larger capacity (300 trees of depth 8, against 60 of depth 6 for the undefended model) to offset the constraint. The gate still reads the port, the protocol and the backward packet and byte totals ($x_0,x_1,x_3,x_5$), on which it is not monotone. Constraints are enforced by the library (XGBoost, LightGBM and CatBoost monotone constraints) and checked by a randomized test that pushes feasible transformations through gates of all libraries and asserts that the margin never falls; a companion test shows that an unconstrained model does fall, so the check is not vacuous.

\begin{proposition}[Conditional certificate]
\label{prop:certificate}
Let $G$ be non-decreasing in $x_2,x_4,x_6,x_{11}$ and independent of $x_7,\dots,x_{10}$, and let $x'$ be derived from a flow $x$ by actions that leave $x_0,x_1,x_3,x_5$ unchanged and do not decrease $x_2,x_4,x_6,x_{11}$ (up to the verifier's absolute tolerance of $10^{-2}$, below the integer resolution of these columns). Then $G(x')\ge G(x)$, so $G(x)\ge\tau$ implies $G(x')\ge\tau$ and the fused score $\max(\text{free},G)$ of $x'$ stays above its alert threshold.
\end{proposition}
The proof is immediate from the definition of monotonicity (Online Resource~1, Section~S1) and does not depend on the generators. The assumptions are strong and we state them explicitly. (i) \emph{The backward features $x_3,x_5$ and the port and protocol do not change.} The gate reads them and is not monotone in them, and chaff, padding, delay and especially re-segmentation can change server responses (more segments can elicit more ACKs), so the certificate does not cover attacks that move the backward direction. Hiding $x_3$ and $x_5$ from the gate (G35) would remove this assumption, but it costs about 4 accuracy and 19 F1 points on CSE-CIC-IDS2018 and gives 36\% benign false positives on UNSW-NB15, and V2 rebuilt on G35 has a whole-cohort evasion of 35--36\% against FFE-v1 and 78--79\% against the white-box attacker on CSE-CIC-IDS2018, against at most 9.0\% for V2 (Online Resource~1, Section~S6). \textbf{V2's guarantee therefore rests on the assumption that backward traffic does not change.} (ii) The attacker cannot remove packets or bytes, split or merge flows, or exploit the feature extractor. (iii) The detector reads the twelve schema features as exported. (iv) The certificate holds if the library's monotone constraint holds, which we verified by randomized test, not by proof. (v) It is relative to the attacker's chosen original flow: a lower-footprint variant of the same attack flow (fewer packets, bytes or less duration) is a different original flow and is not covered, so the certificate is not a worst-case bound on an attack. (vi) It protects the flows the gate flags; it says nothing about flows it does not flag, about the false-positive rate or about poisoning.

\paragraph{Component 2: free classifier and fusion.} Three options were tried on validation: \emph{orbit} training (invariance by augmentation with 100,000 randomly transformed copies of flows of both classes that keep their true labels), a lighter variant, and iterative \emph{adversarial training} on flows crafted by FFE-v1 ($M=5,20$) and the white-box attacker (Online Resource~1, Section~S6). The fused score is $\max(\text{free},\text{gate shifted to fire at }\tau)$. For each library and corpus the pair (free classifier, $\tau\in\{\text{none},0.5,0.8,0.95,0.99\}$) is chosen on the validation split only: among pairs whose validation accuracy, F1 and recall do not fall by more than 0.5 point below the undefended model, we take the lowest mean validation Real ASR of FFE-v1 ($M=20$) and the white-box attacker on a 300-flow cohort. This selection is \textbf{V2}; on four of the nine cells it picks the undefended model as free component, so the gain there comes from the gate alone (Online Resource~1, Table~S11). \textbf{V2-R} allows accuracy and F1 to fall by up to 3 points; it is a post-hoc second tier, defined after a smoke run on UNSW-NB15 LightGBM in which its validation table and the undefended attack rates had been seen (before any full test result), and differs from V2 only on UNSW-NB15; V2 is the headline. The earlier defenses are kept as baselines: \emph{MPS} hides duration and the two rates, and \emph{LGAT} augments training with 5,000 FFE-v1 flows at $M=5$ (Online Resource~1, Algorithm~S4).

\subsection{Protocol}

The cohort is 1,000 random malicious test flows detected by the undefended model of the same library (cohorts differ across libraries); each model is attacked on the cohort flows it still detects, and Real ASR is computed on those. The attackers are FFE-v1 at $M=5$ and $M=20$, CAPGD with a surrogate fitted to the \emph{defended} model's soft outputs, and an adaptive white-box threshold-crossing attacker that reads the defended trees including the gate (a generalization of the control of Section~\ref{sec:experiments} to any library, which also tries one step either side of each threshold). FFE-v2 (with snapping) was run against retrained hardened models only on CIC-IDS2017 XGBoost, first 300 cohort flows: Real ASR / whole-cohort evasion are 100\% / 100\% undefended, 1.0\% / 1.3\% for V2, 8.4\% / 8.7\% for V2 rebuilt on G35 and 0.0\% / 2.3\% for G35 alone; it was not run on any other cell. The gate is covered by Proposition~\ref{prop:certificate} under its assumptions, the free component is not. The free classifiers trained on FFE-v1 and white-box flows are evaluated with the same attack code, so part of the protection is against the training attacks. One seed and one split per cell; Wilson intervals and McNemar tests cover flow sampling only.

\subsection{Results}
\label{sec:defense_results}

\begin{table}[!htbp]
\caption{Defenses (Real ASR, \%, on flows still detected). U: undefended; V2 (bold): gate fused with a free classifier; FFE: FFE-v1, $M=20$; WB: white-box attacker. Acc.\ and Rec.\: clean accuracy and attack recall of U and V2. Real ASR is computed on the flows the model still detects; V2 whole cohort is the maximum whole-cohort evasion over the three attackers (Online Resource~1, Tables~S8--S10, S13). Cohort of 1,000 detected flows per cell; one seed.}\label{tab:defense_evaluation}
\centering
\resizebox{\textwidth}{!}{\begin{varwidth}{\maxdimen}\noindent%
\begin{tabular}{@{}llcccccccccccc@{}}
\toprule
\textbf{Corpus} & \textbf{Library} & \textbf{Acc.\ U/V2} & \textbf{Rec.\ U/V2} & \multicolumn{3}{c}{\textbf{Undefended}} & \multicolumn{3}{c}{\textbf{LGAT}} & \multicolumn{3}{c}{\textbf{V2}} & \textbf{V2 whole} \\
& & & & FFE & CAPGD & WB & FFE & CAPGD & WB & FFE & CAPGD & WB & cohort \\
\midrule
CSE-CIC-IDS2018 & XGBoost & 97.30/97.29 & 77.6/78.4 & 99.5 & 99.0 & 99.5 & 73.1 & 96.1 & 99.5 & \textbf{8.1} & \textbf{0.3} & \textbf{8.1} & 8.3 \\
 & LightGBM & 97.28/97.27 & 77.8/78.2 & 99.8 & 99.5 & 99.8 & 56.8 & 96.3 & 99.8 & \textbf{7.9} & \textbf{0.2} & \textbf{7.9} & 7.9 \\
 & CatBoost & 97.26/97.25 & 77.6/77.9 & 59.8 & 99.3 & 100.0 & 45.6 & 83.8 & 100.0 & \textbf{0.0} & \textbf{0.0} & \textbf{8.8} & 9.0 \\
\midrule
CIC-IDS2017 & XGBoost & 99.96/99.97 & 99.8/99.9 & 99.1 & 77.0 & 100.0 & 31.9 & 52.6 & 100.0 & \textbf{0.7} & \textbf{0.7} & \textbf{0.7} & 0.8 \\
 & LightGBM & 99.96/99.97 & 99.9/99.9 & 34.6 & 57.9 & 35.7 & 33.4 & 36.3 & 35.6 & \textbf{0.4} & \textbf{0.4} & \textbf{0.4} & 0.4 \\
 & CatBoost & 99.85/99.93 & 99.4/99.9 & 32.0 & 5.4 & 33.6 & 31.2 & 31.0 & 33.2 & \textbf{0.1} & \textbf{0.1} & \textbf{0.1} & 0.1 \\
\midrule
UNSW-NB15 & XGBoost & 88.51/88.38 & 88.0/90.4 & 62.2 & 24.4 & 90.9 & 15.0 & 6.2 & 81.7 & \textbf{12.2} & \textbf{1.6} & \textbf{25.4} & 25.5 \\
 & LightGBM & 88.47/88.09 & 87.9/90.5 & 66.4 & 32.7 & 89.5 & 14.5 & 6.3 & 72.3 & \textbf{22.0} & \textbf{2.1} & \textbf{31.9} & 32.4 \\
 & CatBoost & 86.67/86.42 & 89.6/92.6 & 72.6 & 27.5 & 60.1 & 27.1 & 9.0 & 45.4 & \textbf{4.9} & \textbf{1.0} & \textbf{5.7} & 5.7 \\
\bottomrule
\end{tabular}
%
\end{varwidth}}
\end{table}

Table~\ref{tab:defense_evaluation} gives the main result; all models, all attacks and the retained share of the cohort are in Online Resource~1, Tables~S8--S10. The old defenses do not survive the adaptive attackers. MPS and LGAT leave 33--100\% white-box Real ASR on every corpus and library, and LGAT 99.5--100\% on CSE-CIC-IDS2018 and for XGBoost on CIC-IDS2017; its protection against FFE-v1 collapses when the budget moves from the training value $M=5$ to $M=20$ (XGBoost on CSE-CIC-IDS2018: 25.6\% to 73.1\%).

V2 changes this where the certificate applies. On CIC-IDS2017 its Real ASR is at most 0.7\% against every attacker and library, with recall higher than the undefended model by 0.07--0.46 point; the stand-alone gate flags at least 99.8\% of the cohort. The last column of Table~\ref{tab:defense_evaluation} gives V2's whole-cohort evasion (flows the model misses count as evaded; maximum over the three attackers): 0.1--0.8\% on CIC-IDS2017, 0.2--9.0\% on CSE-CIC-IDS2018 and 1.0--32\% on UNSW-NB15, against 24--100\% undefended; the stand-alone gate's whole-cohort evasion equals its miss rate (8--9\%, 0.1--0.2\%, 4.5--6\%). On CSE-CIC-IDS2018 it is 0.0--0.3\% against CAPGD and 0.0--8.1\% against FFE-v1, but 8.8\%, 7.9\% and 8.1\% against the white-box attacker (CatBoost, LightGBM, XGBoost), compared with 99.5--100\% for the undefended models and for LGAT. The remaining 8\% are cohort flows that the gate does not flag and the free classifier loses under attack; V2 keeps at least 91\% of the whole cohort detected under every attack, so the residual is a coverage gap of the gate, not an attack on the certificate. Free classifiers alone stay about 80\% evadable by the white-box attacker on CSE-CIC-IDS2018; the gate closes that gap. Accuracy changes by between $-0.38$ and $+0.08$ point on every cell.

UNSW-NB15, with overlapping classes and about 88\% clean accuracy, is the hard case. At strict parity V2 leaves 12.2\%, 22.0\% and 4.9\% FFE-v1 and 25.4\%, 31.9\% and 5.7\% white-box Real ASR for XGBoost, LightGBM and CatBoost, against 24--91\% undefended; CAPGD reaches 1.0--2.1\%. It is not uniformly better than the older defense: LGAT is clearly better against FFE-v1 for LightGBM (14.5\% against 22.0\% at $M=20$) and indistinguishable for XGBoost at $M=5$ ($p=0.86$), while V2 is better against the white-box attacker and CAPGD. The share of the cohort retained under attack is only 67--94\%. The relaxed tier V2-R, which gives up 2.1--2.4 accuracy points and gains 6.4--7.3 recall points relative to the undefended model, reaches 0.5--4.7\% against every attacker and retains 95--98\%, so on this corpus robustness costs accuracy.

\paragraph{Price of the certificate.} The stand-alone gate loses about 0.7 accuracy, 3.7--3.8 F1 and 5.4--5.8 recall points on CSE-CIC-IDS2018, 1.6--1.8 accuracy points on UNSW-NB15 for XGBoost and LightGBM, and nothing on CIC-IDS2017; it misses 0.1--9\% of the cohort that the undefended model flags (Online Resource~1, Table~S7), and its Real ASR, like every entry of Table~\ref{tab:defense_evaluation}, is computed on the flows still detected, so a 0.0\% for the gate is not complete coverage. An unconstrained 300$\times$8 model with the same hidden columns is as accurate as the undefended model and equally evadable, so the cost is the monotonicity itself, not capacity (0.55 accuracy and 3.1 F1 points on CSE-CIC-IDS2018, 2.3--2.8 accuracy points on UNSW-NB15, 0.06 on CIC-IDS2017). Fusion recovers clean parity. V2's benign false-positive rate is 0.11--0.14\% on CSE-CIC-IDS2018 (undefended 0.03--0.07\%) and 13--19\% on UNSW-NB15 (11--16\%), with no monotone drift found. Building V2 takes 5--33 times the undefended training time and inference runs two ensembles.

  \paragraph{What is not shown.} Random forests (supported by the recipe) and the manipulation-signature detector were not evaluated; the white-box attacker is greedy single-feature threshold crossing, so a stronger search could raise the residual of the free component but not of the gate; validation used a 300-flow cohort and a small grid; and the monotone direction is a modelling assumption that removes the ability to learn ``short and small is malicious'' patterns, which is the source of the accuracy cost. Beyond Proposition~\ref{prop:certificate} we claim no robustness for unflagged flows.

\section{Discussion and Conclusions}
\label{sec:conclusion}

\subsection{What the Results Support, and What They Do Not}
\label{sec:discussion}

\paragraph{The attack.} With score access alone, FFE-v2 matches or exceeds our greedy threshold-crossing control in nine of nine corpus-library cells of a single frozen evaluation (significantly in four) and is above FFE-v1 in seven; the five-seed means point the same way (Tables~\ref{tab:crosslib} and~\ref{tab:seeds}). The gain over FFE-v1 comes from repairing diagnosed weaknesses, in particular the window target and the missing payload-free packet moves, at a price of 28--38 queries (109 on CIC-IDS2017 CatBoost), more than FFE-v1 on the easy corpora and far more than the control (2) or CAPGD (about 3 with an uncounted surrogate). The result is conditional on detection, on the 12-feature schema and on a cost model that FFE steers by. Three limits matter. First, a cell where all methods agree is not an attack success: CatBoost on CIC-IDS2017 (36.4\%) measures how much that model leans on immutable features. Second, the cost story is mixed: on UNSW-NB15 FFE-v2's evasions need more operational cost than FFE-v1's on the flows both evade. Third, the protocol (validation tuning, frozen configuration, one test evaluation) is not a clean hold-out, since the failures that motivated FFE-v2 were seen on the test cohorts, the cohorts are single draws, and the five-seed study, though cleaner, used a configuration tuned on the main split. The claim that FFE-v2 is the strongest attack here is limited to the attacks we ran: the control is our own implementation and is erratic across seeds on CIC-IDS2017 LightGBM (95.7\%, 41.0\%, 95.7\%), CAPGD was run on XGBoost in the single-split test and on three seeds of the other libraries, and LT-Attack and Cascioli et al.\ were not run.

\paragraph{The defense.} The certificate is conditional and narrow. If the port, protocol and backward features stay fixed and the four increase-only columns do not decrease, every flow the monotone gate flags stays flagged; because the gate reads the backward features without being monotone in them, and segmentation or chaff may change server responses, this is not a worst-case bound on an attack, and a lower-footprint variant of the same attack flow is not covered. The stand-alone gate's 0.0\% Real ASR is computed on the flows it still detects and is not a measure of coverage. The residual is about 8\% of the CSE-CIC-IDS2018 cohort for all three libraries and 5--32\% of UNSW-NB15, and the free classifier that covers it carries no guarantee. The gate costs accuracy where ``short and small'' flows are informative (0.7 points on CSE-CIC-IDS2018, 1.6--1.8 on UNSW-NB15), confounded with its larger capacity until a capacity-matched unconstrained arm is run; fusion recovers it on CSE-CIC-IDS2018 but only partly on UNSW-NB15, where LGAT remains better against FFE-v1 for LightGBM. The free classifiers were trained against FFE-v1 and the white-box attacker and are evaluated with the same attack code. Hiding the backward columns from the gate (which would remove the assumption) costs about 4 accuracy and 19 F1 points on CSE-CIC-IDS2018, 36\% benign false positives on UNSW-NB15, and makes V2 much weaker (CSE-CIC-IDS2018 whole-cohort evasion 35--36\% against FFE-v1 and 78--79\% against the white-box attacker, against 8\% for V2), so the guarantee rests on backward traffic not changing. The gate's accuracy cost is the monotonicity, not capacity (an unconstrained 300$\times$8 control is as accurate and equally evadable). FFE-v2 was run against the hardened models only on CIC-IDS2017 XGBoost (V2: 1.0\%, V2 on G35: 8.4\%); V2-R is a post-hoc tier defined after a smoke run on one UNSW-NB15 cell; random forests, a stronger adaptive attacker and more seeds are open.

\paragraph{Feasibility and validity.} ``Feasible'' means conformant to the 12-feature schema verifier and, for FFE-v2, replayed on synthetic traces (Section~\ref{sec:ablation}). The schema holds backward packets fixed, which is exact for pure-ACK chaff but optimistic for re-segmentation; the accounting of chaff bytes decides whether 100\% or 69.2\% of the UNSW-NB15 evasions survive replay; payload semantics were not tested; and a deployed detector may read features outside the twelve. The defense inherits these assumptions. Cohorts contain only flows that the models detect when clean (77.6\% of the CSE-CIC-IDS2018 test attacks), MOEVA and PSO ran on 200 flows, CAPGD and the NIDSGAN-style generator are our re-implementations with one small surrogate, and the random split lets near-duplicates fall on both sides so accuracies are optimistic.

\subsection{Concluding Remarks and Future Directions}
\label{sec:concluding_remarks}

This article measured, at full scale and across three tree libraries, how far a schema-constrained attacker with score access can evade flow-based NIDS, and what a defender can promise in return. A diagnosis of FFE-v1 on validation data led to FFE-v2: a re-segmentation generator that preserves the Lie-cone invariants, multi-scale probing, window targets from the attacker's benign sample, log-odds progress and cost refinement. In one frozen evaluation (not a clean hold-out) FFE-v2 reaches 99.5\%, 100.0\% and 98.3\% Real ASR on XGBoost, equals or exceeds our greedy white-box control in all nine cells and repairs the LightGBM and CatBoost failures of FFE-v1, with 28--38 queries and a higher operational cost than FFE-v1 on UNSW-NB15; CatBoost on CIC-IDS2017 is capped at 36.4\% by immutable features. OR-fusing a monotone gate with a free classifier (hardened on five of nine cells) cuts Real ASR on the flows still detected, against FFE-v1, CAPGD and a greedy white-box attacker, from 5--100\% to at most 8.8\% on CSE-CIC-IDS2018 and CIC-IDS2017 and to 1.0--32\% on UNSW-NB15, with a conditional guarantee for flagged flows. Future work includes FFE-v2 against the hardened models, replay on a real network stack including server ACK behavior, the official white-box tree attacks, random forests and neural detectors, time- or host-grouped splits, and further public corpora.


\backmatter

\bmhead{Supplementary information}
The accompanying file Online Resource~1 contains the proofs, auxiliary algorithms and additional experiments referred to in the text (Sections~S1--S13, Tables~S1--S21, Algorithms~S1--S5).

\section*{Declarations}

\bmhead{Funding}
This research received no external funding. [To be confirmed by the author.]

\bmhead{Competing interests}
The author declares no known competing financial interests or personal relationships that could have appeared to influence the work reported in this paper.

\bmhead{Ethics approval}
Not applicable. The study uses only public network-flow datasets and involves no human participants or animals. No real network was attacked.

\bmhead{Consent to participate}
Not applicable.

\bmhead{Consent for publication}
Not applicable.

\bmhead{Data availability}
All flow data come from the public CSE-CIC-IDS2018 and CIC-IDS2017 archives of the Canadian Institute for Cybersecurity (University of New Brunswick, \url{https://www.unb.ca/cic/datasets/}) and the UNSW-NB15 archive of UNSW Canberra (\url{https://research.unsw.edu.au/projects/unsw-nb15-dataset}). The dataset builders and the per-flow result archives for FFE-v1, FFE-v2, the white-box control, all baselines, the validation tuning, the multi-seed runs, the packet replay and the defenses are in the repository \url{https://github.com/vinhqdang/adversarial_ids}.

\bmhead{Code availability}
The code of all attacks, baselines and defenses is at \url{https://github.com/vinhqdang/adversarial_ids}; every table is regenerated from the result files by \texttt{experiments/summarize\_results.py}.

\bmhead{Author contributions}
Quang-Vinh Dang: Conceptualization, Methodology, Software, Formal analysis, Investigation, Writing -- original draft, Writing -- review \& editing.

\bmhead{Use of AI tools}
Generative AI tools were used under the author's supervision to assist with software development, experiment scripting and drafting and editing of the text; the author reviewed and verified all code, results and text and takes full responsibility for the content.

\bmhead{Dual use and responsible disclosure}
This article studies the evasion of intrusion detectors. The attacks are red-teaming tools that estimate what a schema-constrained adversary with score access can reach, and no real network was attacked; the discussion frames the results as a robustness audit and the monotone-gate defense of Section~\ref{sec:defense} is offered together with its limits so that defenders can act on the findings.

\bibliography{references}

\end{document}
